% Ch2 self-study -- "I do / You do" ladder for Zumdahl 2.1-2.8, four
% YOUR TURN questions per worked example. Isotopes and compounds chosen
% disjoint from ch02-notes, ws1, ws2 and especially ws3 (which owns the
% nomenclature drilling -- ladder 7 here teaches the DECISION, not the
% vocabulary, and sends the reader to ws3 for volume).
\documentclass{apchem-worksheet}

\apcunitword{Chapter}
\apcunit{2}
\apcunittitle{Atoms, Molecules, and Ions}
\apcdoctitle{Self-Study \textbullet\ Chapter 2, I~do / You~do}
\apcsubtitle{Zumdahl \S2.1--2.8 \textbullet\ PDF pp.\ 70--95 \textbullet\ four YOUR TURN questions per ladder}

\begin{document}
\apctitleblock

\begin{apcnote}[How to use these notes --- read this first]
Each skill is a \textbf{ladder}: a fully worked example in the
solid-framed box, then \textbf{four} YOUR TURN questions in the dashed
box --- same skill, new numbers, no help.

\begin{enumerate}[leftmargin=*,itemsep=0.15em]
  \item \textbf{Work all four} before checking anything.
  \item Check against the gray \emph{check:} line. All four right on the
        first try earns the tick on the tracker (last page).
  \item Any misses: re-read the worked example, then redo the missed ones
        from a blank page.
\end{enumerate}

\textbf{What this chapter is really about.} Sections 2.1--2.4 are the
\emph{evidence}: a century of experiments, each forcing one conclusion
about what atoms must be like. Sections 2.5--2.8 are the \emph{result}:
how we count, charge, and name them. Learn the evidence as
experiment~$\to$~conclusion pairs and the history stops being a list of
names to memorize.
\end{apcnote}

% =====================================================================
\section*{\sffamily Ladder 1 \textbullet\ The laws that demanded atoms}
\zsec{2.2--2.3}

Three quantitative laws came first; Dalton's theory was invented to
explain them.

\begin{center}
\renewcommand{\arraystretch}{1.25}
\begin{tabular}{@{}p{3.6cm}p{9.6cm}@{}}
\hline
\textbf{Law} & \textbf{What it states} \\
\hline
Conservation of mass & mass is neither created nor destroyed in a
  reaction \\
Definite proportion & a given compound always has the same proportions
  by mass \\
Multiple proportions & when two elements form more than one compound,
  the masses of one combining with a fixed mass of the other are in
  \emph{small whole-number ratios} \\
\hline
\end{tabular}
\end{center}

The third is the decisive one: small whole numbers are what you expect if
matter comes in indivisible units, and are hard to explain otherwise.

\begin{workedexample}[I do: testing the law of multiple proportions]
Two oxides of sulfur are analysed. Oxide A is 50.05\% S by mass; oxide B
is 40.05\% S. Show that these obey the law of multiple proportions, and
identify the ratio.

\textbf{Fix the sulfur, compare the oxygen.} Take \SI{100}{\gram} of
each:
\[ \text{A}: 50.05~\text{g S} + 49.95~\text{g O} \qquad
   \text{B}: 40.05~\text{g S} + 59.95~\text{g O} \]

\textbf{Oxygen per gram of sulfur:}
\[ \text{A}: \frac{49.95}{50.05} = 0.998 \qquad
   \text{B}: \frac{59.95}{40.05} = 1.497 \]

\textbf{Ratio of those two figures:}
\[ \frac{1.497}{0.998} = 1.50 = \frac{3}{2} \]

A small whole-number ratio, exactly as the law requires. The compounds
are \ce{SO2} and \ce{SO3}: for a fixed amount of sulfur, B holds half
again as much oxygen.

\textbf{The move to remember:} hold one element constant, then divide.
The ratio only looks simple after that normalization.
\end{workedexample}

\begin{yourturn}[YOUR TURN 1 --- four questions]
\begin{enumerate}[leftmargin=*,itemsep=0.5em]
  \item Two carbon oxides give \SI{1.33}{\gram} and \SI{2.66}{\gram} of
        oxygen per gram of carbon. State the ratio and name both
        compounds. \workspace[1.6cm]
  \item A \SI{10.0}{\gram} sample of a compound decomposes into
        \SI{4.0}{\gram} of A and \SI{6.0}{\gram} of B. What mass of B is
        in a \SI{25.0}{\gram} sample, and which law did you use?
        \workspace[1.8cm]
  \item \SI{12.0}{\gram} of carbon burns in \SI{32.0}{\gram} of oxygen to
        give carbon dioxide. What mass of \ce{CO2} forms, and by which
        law? \workspace[1.6cm]
  \item Which of the three laws is the strongest evidence \emph{for}
        atoms, and why? \workspace[1.6cm]
\end{enumerate}
\selfcheck{(a) $1{:}2$ --- \ce{CO} and \ce{CO2} \quad (b)
\SI{15.0}{\gram} of B, by definite proportion \quad (c)
\SI{44.0}{\gram}, by conservation of mass \quad (d) multiple
proportions}
\keyonly{\begin{apcanswer}(a) $2.66/1.33 = 2.00$, so \textbf{1:2} ---
\ce{CO} and \ce{CO2}. (b) The compound is 60\% B by mass in any sample:
$25.0 \times 0.600 = \mathbf{15.0}$~g, by the law of \textbf{definite
proportion}. (c) $12.0 + 32.0 = \mathbf{44.0}$~g, by \textbf{conservation
of mass} --- nothing is lost. (d) \textbf{Multiple proportions.} Definite
proportion says a compound is consistent; multiple proportions says the
\emph{alternatives} differ by small whole numbers, which is what you
expect if atoms combine in fixed counts and is awkward to explain if
matter is continuous.\end{apcanswer}}
\end{yourturn}

% =====================================================================
\section*{\sffamily Ladder 2 \textbullet\ Three experiments, three
conclusions}
\zsec{2.4}

Every one of these is examinable as \emph{experiment $\to$ observation
$\to$ conclusion}. Learn the chain, not the anecdote.

\begin{center}
\renewcommand{\arraystretch}{1.25}
\begin{tabular}{@{}p{2.9cm}p{4.6cm}p{5.6cm}@{}}
\hline
\textbf{Experiment} & \textbf{Key observation} & \textbf{Conclusion} \\
\hline
Thomson, cathode ray & the beam bends toward the positive plate,
  identically for every cathode metal &
  a negative particle common to all atoms; gave $e/m$, not $e$ \\
Millikan, oil drop & droplet charges are all multiples of one value &
  the charge on the electron, so with Thomson's ratio, its mass \\
Rutherford, gold foil & a few alpha particles deflect sharply backward &
  a tiny, dense, positive nucleus in a mostly empty atom \\
\hline
\end{tabular}
\end{center}

\begin{workedexample}[I do: reasoning from the observation]
\textbf{Why did Rutherford's result overturn the plum-pudding model?}

\textbf{What the old model predicted.} If positive charge were spread
thinly through the whole atom, an alpha particle would meet only weak,
diffuse repulsion everywhere. \emph{Every} alpha should pass nearly
straight through, with slight deflection at most.

\textbf{What was observed.} The vast majority did pass straight through
--- but roughly 1 in 20\,000 bounced back at more than
\SI{90}{\degree}.

\textbf{Why that is decisive.} Reversing a fast, massive alpha particle
requires a concentrated positive charge that is also \emph{massive}
enough not to simply recoil away. Spread-out charge cannot do it at any
strength. So the positive charge and nearly all the mass must sit in a
minute volume --- a \term{nucleus} --- while the rest of the atom is
essentially empty.

\textbf{Note the shape of the argument:} the rare event, not the common
one, carries the information. Most alphas passing through told
Rutherford the atom is mostly empty; the few that came back told him
what sits at the centre.
\end{workedexample}

\begin{yourturn}[YOUR TURN 2 --- four questions]
\begin{enumerate}[leftmargin=*,itemsep=0.5em]
  \item Thomson used several different cathode metals and got the same
        particle each time. What does that show? \workspace[1.6cm]
  \item Millikan measured charge. Why was Thomson's experiment alone not
        enough to give the electron's mass? \workspace[1.6cm]
  \item In the gold foil experiment, what did the alphas that passed
        straight through establish? \workspace[1.6cm]
  \item Which subatomic particle was discovered last, and why was it the
        hardest to detect? \workspace[1.6cm]
\end{enumerate}
\selfcheck{(a) electrons are in all atoms \quad (b) it gave only the
charge-to-mass ratio \quad (c) the atom is mostly empty space \quad
(d) the neutron --- it carries no charge}
\keyonly{\begin{apcanswer}(a) That \textbf{electrons are a component of
all atoms}, not something peculiar to one element --- which is why the
atom could no longer be considered indivisible. (b) Thomson measured only
the \textbf{charge-to-mass ratio} $e/m$. Neither quantity alone is
obtainable from a ratio; Millikan's charge measurement supplied the
second equation, and the mass followed. (c) That the atom is
\textbf{mostly empty space} --- the nucleus occupies a vanishing fraction
of the volume. (d) The \textbf{neutron} (Chadwick, 1932). Being
\textbf{electrically neutral} it is unaffected by electric and magnetic
fields, so none of the deflection techniques that revealed the electron
and the proton could detect it.\end{apcanswer}}
\end{yourturn}

% =====================================================================
\section*{\sffamily Ladder 3 \textbullet\ Reading isotope symbols}
\zsec{2.5}

\[ \ce{^{A}_{Z}X} \qquad
   Z = \text{atomic number} = \text{protons} \qquad
   A = \text{mass number} = \text{protons} + \text{neutrons} \]

Three counts, three rules:
\begin{itemize}[leftmargin=*,itemsep=0.15em]
  \item \textbf{Protons} $= Z$ --- and $Z$ \emph{is} the element's
        identity. Change it and you have a different element.
  \item \textbf{Neutrons} $= A - Z$. Change it and you have an
        \term{isotope} of the same element.
  \item \textbf{Electrons} $= Z - \text{charge}$. Change it and you have
        an ion.
\end{itemize}

\begin{workedexample}[I do: three species, one procedure]
\textbf{(a) \ce{^{31}_{15}P}:} protons $= 15$; neutrons $= 31 - 15 =
\mathbf{16}$; neutral, so electrons $= \mathbf{15}$.

\textbf{(b) \ce{^{80}_{35}Br^-}:} protons $= 35$; neutrons $= 80 - 35 =
\mathbf{45}$; the $1-$ charge means one \emph{extra} electron, so
electrons $= 35 - (-1) = \mathbf{36}$.

\textbf{(c) \ce{^{27}_{13}Al^{3+}}:} protons $= 13$; neutrons
$= 27 - 13 = \mathbf{14}$; the $3+$ charge means three electrons
\emph{missing}, so electrons $= 13 - 3 = \mathbf{10}$.

\textbf{The sign trap, stated once:} a \emph{negative} ion has
\emph{more} electrons, a \emph{positive} ion has fewer. The formula
electrons $= Z - \text{charge}$ handles both automatically, provided you
carry the charge with its sign.

\textbf{Worth noticing:} the bromide ion in (b) has 36 electrons, the
same as krypton, and the aluminium ion in (c) has 10, the same as neon.
Ions form by reaching a noble-gas electron count.
\end{workedexample}

\begin{yourturn}[YOUR TURN 3 --- four questions]
\begin{enumerate}[leftmargin=*,itemsep=0.5em]
  \item \ce{^{56}_{26}Fe}: protons, neutrons, electrons.
        \workspace[1.4cm]
  \item \ce{^{32}_{16}S^{2-}}: protons, neutrons, electrons.
        \workspace[1.4cm]
  \item An ion has 20 protons, 20 neutrons, and 18 electrons. Identify it
        and write its full symbol. \workspace[1.6cm]
  \item \ce{^{14}C} and \ce{^{14}N} share a mass number. Are they
        isotopes of each other? Explain. \workspace[1.6cm]
\end{enumerate}
\selfcheck{(a) 26, 30, 26 \quad (b) 16, 16, 18 \quad
(c) \ce{^{40}_{20}Ca^{2+}} \quad (d) no --- different $Z$, so different
elements}
\keyonly{\begin{apcanswer}(a) $p = 26$, $n = 56 - 26 = \mathbf{30}$,
$e = \mathbf{26}$ (neutral). (b) $p = 16$, $n = \mathbf{16}$,
$e = 16 - (-2) = \mathbf{18}$. (c) $Z = 20$ is calcium;
$A = 20 + 20 = 40$; two electrons missing means a $2+$ charge:
\textbf{\ce{^{40}_{20}Ca^{2+}}}. (d) \textbf{No.} Isotopes share $Z$ and
differ in $A$. These share $A$ but have $Z = 6$ and $Z = 7$ --- different
elements entirely. Species with equal $A$ are called isobars, which is a
different relationship.\end{apcanswer}}
\end{yourturn}

% =====================================================================
\section*{\sffamily Ladder 4 \textbullet\ The periodic table as a map}
\zsec{2.7}

The table is arranged so that position predicts behaviour.

\begin{itemize}[leftmargin=*,itemsep=0.15em]
  \item \textbf{Groups} (columns) share chemical behaviour: 1A
        \term{alkali metals}, 2A \term{alkaline earth metals}, 7A
        \term{halogens}, 8A \term{noble gases}.
  \item \textbf{Periods} (rows) are filled left to right.
  \item \textbf{Metals} sit left of the staircase, \textbf{nonmetals}
        right, \term{metalloids} along it (B, Si, Ge, As, Sb, Te).
  \item Metals \emph{lose} electrons to form cations; nonmetals
        \emph{gain} them to form anions.
\end{itemize}

\begin{workedexample}[I do: predicting from position alone]
\textbf{Predict the ion formed by strontium, by selenium, and by argon
--- using only their positions.}

\textbf{Strontium} is in group 2A, so it has two outer electrons to lose
and forms \textbf{\ce{Sr^2+}}. Every group 2A metal does the same, which
is why the group has a shared name.

\textbf{Selenium} is in group 6A, two short of eight, so it gains two:
\textbf{\ce{Se^2-}}. The pattern across the nonmetals is
$8 - \text{group number}$ electrons gained.

\textbf{Argon} is in group 8A, already at eight outer electrons, so it
forms \textbf{no ion} --- there is nothing to gain or lose cheaply. That
is exactly why the noble gases are unreactive, and why every other ion
above lands on a noble-gas count.

\textbf{The single rule underneath all three:} elements form the ion that
reaches the nearest noble-gas electron count.
\end{workedexample}

\begin{yourturn}[YOUR TURN 4 --- four questions]
\begin{enumerate}[leftmargin=*,itemsep=0.5em]
  \item Predict the ions formed by potassium, by nitrogen, and by
        barium. \workspace[1.6cm]
  \item Which group do \ce{Cl}, \ce{Br}, and \ce{I} belong to, and what
        charge do their ions carry? \workspace[1.4cm]
  \item Name the metalloid in period 3, and say what ``metalloid'' means.
        \workspace[1.6cm]
  \item Predict the formula of the compound formed between magnesium and
        nitrogen, showing the charge balance. \workspace[1.8cm]
\end{enumerate}
\selfcheck{(a) \ce{K+}, \ce{N^3-}, \ce{Ba^2+} \quad (b) group 7A, the
halogens; $1-$ \quad (c) silicon --- properties between metal and
nonmetal \quad (d) \ce{Mg3N2}}
\keyonly{\begin{apcanswer}(a) \textbf{\ce{K+}} (1A loses one),
\textbf{\ce{N^3-}} (5A gains three), \textbf{\ce{Ba^2+}} (2A loses two).
(b) Group \textbf{7A}, the \textbf{halogens}; each gains one electron to
form a $\mathbf{1-}$ ion. (c) \textbf{Silicon} --- a metalloid has
properties intermediate between metals and nonmetals, notably
semiconduction, which is why silicon underlies every microchip.
(d) \ce{Mg^2+} with \ce{N^3-}: three $2+$ against two $3-$ gives $+6$ and
$-6$, so \textbf{\ce{Mg3N2}}.\end{apcanswer}}
\end{yourturn}

% =====================================================================
\section*{\sffamily Ladder 5 \textbullet\ Molecules, ions, and formulas}
\zsec{2.6}

An \textbf{ionic} compound is a lattice of ions held by electrostatic
attraction --- its formula is the \emph{simplest whole-number ratio}, not
a discrete particle. A \textbf{molecular} compound is made of discrete
molecules, and its formula counts the atoms in one of them.

Charge balance writes every ionic formula: total positive $=$ total
negative, with no charges shown in the finished formula.

\begin{workedexample}[I do: building formulas from charges]
\textbf{(a) Aluminium and oxygen.} \ce{Al^3+} and \ce{O^2-}. The lowest
common multiple of 3 and 2 is 6, so two \ce{Al^3+} ($+6$) balance three
\ce{O^2-} ($-6$): \textbf{\ce{Al2O3}}.

\textbf{(b) Calcium and phosphate.} \ce{Ca^2+} and \ce{PO4^3-}. Again
LCM 6: three \ce{Ca^2+} ($+6$) with two \ce{PO4^3-} ($-6$). The
polyatomic ion is taken more than once, so it needs parentheses:
\textbf{\ce{Ca3(PO4)2}}.

\textbf{(c) Ammonium and sulfate.} \ce{NH4+} and \ce{SO4^2-}: two
ammonium per sulfate, \textbf{\ce{(NH4)2SO4}}.

\textbf{Why parentheses are not optional:} \ce{NH42SO4} would read as
forty-two hydrogens. The parentheses say ``two of the whole group''.

\textbf{And the conceptual point:} \ce{Al2O3} does not describe a
molecule. There is no \ce{Al2O3} particle --- only a vast lattice in
which aluminium and oxide ions occur in a 2:3 ratio.
\end{workedexample}

\begin{yourturn}[YOUR TURN 5 --- four questions]
\begin{enumerate}[leftmargin=*,itemsep=0.5em]
  \item Write the formula from lithium and sulfide ions.
        \workspace[1.2cm]
  \item Write the formula from iron(III) and hydroxide ions.
        \workspace[1.4cm]
  \item Write the formula from ammonium and phosphate ions.
        \workspace[1.4cm]
  \item Is \ce{P4O10} ionic or molecular? How can you tell from the
        elements alone? \workspace[1.6cm]
\end{enumerate}
\selfcheck{(a) \ce{Li2S} \quad (b) \ce{Fe(OH)3} \quad
(c) \ce{(NH4)3PO4} \quad (d) molecular --- two nonmetals}
\keyonly{\begin{apcanswer}(a) \ce{Li+} and \ce{S^2-}: two lithium per
sulfide, \textbf{\ce{Li2S}}. (b) \ce{Fe^3+} and \ce{OH-}: three
hydroxides, and because the polyatomic ion repeats it takes parentheses,
\textbf{\ce{Fe(OH)3}}. (c) \ce{NH4+} and \ce{PO4^3-}: three ammonium per
phosphate, \textbf{\ce{(NH4)3PO4}}. (d) \textbf{Molecular} ---
phosphorus and oxygen are both nonmetals, and ionic bonding needs a metal
(or ammonium) to donate electrons. Note also that its subscripts, 4 and
10, are not a simplest ratio: a molecular formula counts the atoms in one
molecule rather than reducing.\end{apcanswer}}
\end{yourturn}

% =====================================================================
\section*{\sffamily Ladder 6 \textbullet\ Averaging isotopes}
\zsec{2.5}

The mass on the periodic table is a \textbf{weighted average} over the
naturally occurring isotopes:
\[ \bar m = \sum(\text{fractional abundance} \times \text{isotope mass}) \]
No single atom has that mass. The average always lies between the
extremes and leans toward the abundant isotope --- two checks that cost
nothing and catch most errors.

\begin{workedexample}[I do: the atomic mass of chlorine]
Chlorine is 75.78\% \ce{^{35}Cl} (34.969 u) and 24.22\% \ce{^{37}Cl}
(36.966 u).

\[ \bar m = (0.7578)(34.969) + (0.2422)(36.966) \]
\[ \phantom{\bar m} = 26.500 + 8.953 = \mathbf{35.45~u} \]

\textbf{Both checks:} 35.45 lies between 34.969 and 36.966 --- good; and
it sits much closer to 34.969, which is right because \ce{^{35}Cl} is
roughly three times as abundant.

\textbf{What the value means physically:} there is no chlorine atom
weighing 35.45 u. Every one weighs about 35 or about 37; 35.45 is what a
mole of natural chlorine averages to, which is exactly the number needed
for weighing bulk samples.
\end{workedexample}

\begin{yourturn}[YOUR TURN 6 --- four questions]
\begin{enumerate}[leftmargin=*,itemsep=0.5em]
  \item Copper is 69.17\% \ce{^{63}Cu} (62.930 u) and 30.83\%
        \ce{^{65}Cu} (64.928 u). Find the atomic mass.
        \workspace[1.8cm]
  \item Without calculating: is boron's atomic mass nearer 10 or 11,
        given that \ce{^{11}B} is about 80\% abundant?
        \workspace[1.4cm]
  \item Lithium's atomic mass is 6.94 u, with isotopes of 6.015 and
        7.016 u. Which is more abundant, and how do you know?
        \workspace[1.6cm]
  \item Explain in one sentence why no individual atom has the mass
        printed on the periodic table. \workspace[1.4cm]
\end{enumerate}
\selfcheck{(a) 63.55 u \quad (b) nearer 11 \quad (c) \ce{^{7}Li} --- the
average sits close to 7.016 \quad (d) it is an average over isotopes}
\keyonly{\begin{apcanswer}(a) $(0.6917)(62.930) + (0.3083)(64.928) =
43.529 + 20.017 = \mathbf{63.55}$~u. (b) \textbf{Nearer 11} --- the
average always leans toward the abundant isotope. (c)
\textbf{\ce{^{7}Li}}: 6.94 lies much closer to 7.016 than to 6.015, and
the average leans toward whichever isotope dominates. (d) Because it is a
\textbf{weighted average} over the naturally occurring isotopes; each
individual atom instead has one of the specific isotopic
masses.\end{apcanswer}}
\end{yourturn}

% =====================================================================
\section*{\sffamily Ladder 7 \textbullet\ Choosing the naming rule}
\zsec{2.8}

Nomenclature is drilled at length in Worksheet~3. What matters \emph{here}
is the decision that comes first --- almost every naming error is a
mis-classification, not a forgotten name.

\begin{center}
\renewcommand{\arraystretch}{1.25}
\begin{tabular}{@{}p{1.1cm}p{4.6cm}p{7.4cm}@{}}
\hline
\textbf{Type} & \textbf{How you recognize it} & \textbf{How it is named} \\
\hline
I   & metal of fixed charge (1A, 2A, 3A, \ce{Ag}, \ce{Zn}, \ce{Cd}) &
  metal $+$ anion; \emph{no} Roman numeral \\
II  & transition metal, or \ce{Sn} / \ce{Pb} &
  Roman numeral gives the charge, deduced from the anions \\
III & two nonmetals &
  prefixes on both elements (no \emph{mono-} on the first) \\
A   & formula starts with \ce{H}, aqueous &
  \emph{-ide} to hydro\ldots ic; \emph{-ite} to \emph{-ous};
  \emph{-ate} to \emph{-ic} \\
\hline
\end{tabular}
\end{center}

\begin{workedexample}[I do: classify first, then name]
Name \ce{TiO2}, \ce{SO3}, and \ce{HBrO3(aq)} --- classifying each before
writing anything.

\textbf{\ce{TiO2}.} Titanium is a transition metal $\Rightarrow$
\textbf{Type II}, so a Roman numeral is required and no prefixes are
used. Two \ce{O^2-} give $-4$, so titanium is $+4$:
\textbf{titanium(IV) oxide}. The numeral came from the \emph{anions},
never from the subscript on titanium.

\textbf{\ce{SO3}.} Sulfur and oxygen are both nonmetals $\Rightarrow$
\textbf{Type III}, so prefixes and no numeral: \textbf{sulfur trioxide}.
Writing ``sulfur(VI) oxide'' applies the wrong rule set entirely.

\textbf{\ce{HBrO3(aq)}.} Starts with \ce{H} and is aqueous $\Rightarrow$
\textbf{acid}. The anion \ce{BrO3-} is bromate, an \emph{-ate}, which
gives an \emph{-ic} acid: \textbf{bromic acid}.

\textbf{The pattern:} each name followed from the classification, and
each classification took one glance at what kind of elements were
present.
\end{workedexample}

\begin{yourturn}[YOUR TURN 7 --- four questions]
Classify (I, II, III, or A), then name:
\begin{enumerate}[leftmargin=*,itemsep=0.5em]
  \item \ce{SrBr2} \workspace[1.2cm]
  \item \ce{PCl3} \workspace[1.2cm]
  \item \ce{CoO} \workspace[1.4cm]
  \item \ce{H2SO3(aq)} \workspace[1.4cm]
\end{enumerate}
\selfcheck{(a) I --- strontium bromide \quad (b) III --- phosphorus
trichloride \quad (c) II --- cobalt(II) oxide \quad (d) A --- sulfurous
acid}
\keyonly{\begin{apcanswer}(a) \textbf{Type I} (strontium is 2A, fixed
charge): \textbf{strontium bromide}, no numeral. (b) \textbf{Type III}
(two nonmetals): \textbf{phosphorus trichloride} --- prefixes, and no
\emph{mono-} on the first element. (c) \textbf{Type II} (cobalt is a
transition metal): one \ce{O^2-} forces \ce{Co^2+}, so
\textbf{cobalt(II) oxide}. (d) \textbf{Acid}: the anion is sulfite, an
\emph{-ite}, which gives an \emph{-ous} acid --- \textbf{sulfurous acid},
not sulfuric.\end{apcanswer}}
\end{yourturn}

\begin{apcnote}[Where to get volume on naming]
If Ladder~7 felt shaky, the fix is repetition rather than more theory.
\textbf{Worksheet~3} of this chapter builds the ion tables from memory and
then works 94 items --- type identification, naming by type, formula
writing, extending the oxyanion pattern, and an error hunt. Do it once
untimed, then again timed.
\end{apcnote}

% =====================================================================
\section*{\sffamily Mastery tracker}

Tick a row only if \textbf{all four} YOUR TURN questions were right on
the first attempt.

\renewcommand{\arraystretch}{1.35}
\noindent\begin{tabular}{@{}c p{6.4cm} c p{4.6cm}@{}}
\toprule
\textbf{First try?} & \textbf{Skill} & \textbf{Ladder} &
  \textbf{If not, re-read\ldots}\\
\midrule
$\square$ & the three fundamental laws & 1 & hold one element constant \\
$\square$ & the three classic experiments & 2 & observation to
  conclusion \\
$\square$ & isotope symbols and counts & 3 & electrons $= Z -$ charge \\
$\square$ & the periodic table as a map & 4 & nearest noble-gas count \\
$\square$ & formulas from charge balance & 5 & LCM, then parentheses \\
$\square$ & weighted atomic mass & 6 & between, and leaning \\
$\square$ & choosing the naming rule & 7 & classify before naming \\
\bottomrule
\end{tabular}

\begin{apcnote}[Scoring yourself honestly]
7/7: on to Chapter 3, where every calculation assumes Ladders 3--6.
5--6: redo the missed ladders tomorrow, not today. 4 or fewer: check
whether the misses cluster in the \emph{history} ladders (1--2) or the
\emph{counting} ladders (3--6). History misses are a memorization problem,
fixed by rehearsing the experiment-to-conclusion chains aloud. Counting
misses matter more, because Chapter 3 is built directly on them --- fix
those first.
\end{apcnote}

\end{document}
